\subsection{Zeigerdiagramme 1}
\Aufgabe{

Gegeben sei die Spannung $\hat{U} = 325\ V \cdot \mathrm{e}^{\mathrm{j}30^\circ}$.

\begin{itemize}
    \item [\bf a)] Zeichnen Sie das Zeigerdiagramm für die Spannung $\hat{U}$ im komplexen Zahlenraum.
    \item [\bf b)] Berechnen Sie die Real- und Imaginärteile der komplexen Spannung $\hat{U}$.
    \item [\bf c)] Erklären Sie die Bedeutung des Phasenwinkels in Bezug auf die Wechselspannung.
\end{itemize}
}

\Loesung{
    \begin{itemize}

        \item[\bf a)]
              Zeigerdiagramm für die Spannung $\hat{U}$ \\
              
              \begin{figure}[H]
                  \centering
                  \resizebox{0.3\textwidth}{!}{
                      \input{Tikz/src/Aufgaben/Loesung_Zeigerdiagramme1}\alttext{Das Bild zeigt ein komplexes Koordinatensystem mit Realachse (Re) und Imaginärachse (Im). Vom Ursprung aus ist ein Spannungszeiger Û eingezeichnet, der einen Winkel von 30 Grad zur positiven Realachse bildet. Der Nullwinkel 0° ist entlang der Realachse markiert. }
                  }
              \end{figure}
              
        \item[\bf b)] 
              Die Real- und Imaginärteile der Spannung $\hat{U}$ können mit den folgenden Formeln berechnet werden:
              
              \begin{eqa}
                  U_{\text{Re}} &= \hat{U} \cdot \cos(\varphi) = 325 \cdot \cos(30^\circ) \approx 281,46\ V      \nonumber   \\
                  U_{\text{Im}} &= \hat{U} \cdot \sin(\varphi) = 325 \cdot \sin(30^\circ) = 162,5\ V   \nonumber
              \end{eqa}
              
        \item[\bf c)] 
              Der Phasenwinkel $\varphi$ beschreibt den zeitlichen Versatz zwischen der Spannung und dem Strom in einem Wechselstromkreis. 
              Ein positiver Phasenwinkel deutet darauf hin, dass die Spannung der Strom nachfolgt (induktive Last), während ein negativer 
              Phasenwinkel anzeigt, dass der Strom der Spannung nachfolgt (kapazitive Last). 
              
    \end{itemize}
    
}